\section{Budget et coûts}

\subsection{Méthode}

Le budget du projet est présenté selon deux notions distinctes, conformément
à la grille d'estimation fournie au lancement du projet :

\begin{itemize}
  \item l'\textbf{estimation du coût du prototype} : ce que coûterait la
        réalisation si l'ensemble du travail était facturé aux tarifs
        horaires d'un bureau d'ingénierie (main d'\oe uvre, machines de
        fabrication et pièces achetées) ;
  \item le \textbf{coût réel pour l'école} : la dépense effectivement engagée.
        Les heures de travail des étudiants ne sont pas facturées, et les
        moyens de fabrication (impression 3D, découpe laser) ainsi qu'une
        partie des composants proviennent du stock du MakeLab. Le coût réel
        se résume donc aux pièces effectivement commandées, justifiées par
        les quittances.
\end{itemize}

Les tarifs horaires retenus sont ceux de la grille du cours : ingénieur HES
à 150\,CHF/h, mécanicien à 80\,CHF/h, impression 3D à 40\,CHF/h et découpe
laser à 50\,CHF/h.

\subsection{Estimation du coût du prototype}

Le temps de travail provient du planning final (\autoref{annexe:planning}) :
\textbf{361 heures} consommées au total. La répartition retenue distingue
les heures d'ingénierie (logiciel, électronique, électrique, gestion de
projet) des heures de mécanique (conception, fabrication et montage,
réalisées par le pôle ROCM). Les durées d'impression 3D et de découpe laser
correspondent au temps machine, estimé séparément.

\begin{table}[H]
  \centering
  \begin{tabular}{@{}p{5cm}r r r@{}}
    \toprule
    \textbf{Poste} & \textbf{Quantité} & \textbf{Tarif} & \textbf{Total} \\
    \midrule
    Ingénieur HES        & 314\,h & 150\,CHF/h & 47\,100\,CHF \\
    Mécanicien           & 47\,h  & 80\,CHF/h  & 3\,760\,CHF \\
    Impression 3D        & 15\,h  & 40\,CHF/h  & 600\,CHF \\
    Découpe laser        & 2\,h   & 50\,CHF/h  & 100\,CHF \\
    Pièces achetées      & ---    & ---        & $\approx$\,500\,CHF \\
    \midrule
    \textbf{Total estimé} & & & $\approx$\,\textbf{52\,060\,CHF} \\
    \bottomrule
  \end{tabular}
  \caption{Estimation du coût de développement du prototype, tarifs HES.}
  \label{tab:budget-estimation}
\end{table}

L'essentiel du coût estimé provient de la main d'\oe uvre d'ingénierie : le
projet est avant tout un travail de conception et de développement
logiciel. Le coût matériel reste marginal devant le coût en temps.

\subsection{Pièces achetées — nomenclature}

La nomenclature ci-dessous distingue les composants \emph{commandés} pour le
projet (justifiés par les quittances en annexe) des composants \emph{issus
du stock} du MakeLab, dont le prix est estimé au prix de marché pour
l'estimation du prototype.

\begin{table}[H]
  \centering
  \begin{tabular}{@{}p{6.5cm}p{2.2cm}r p{2.3cm}@{}}
    \toprule
    \textbf{Composant} & \textbf{Origine} & \textbf{Prix} & \textbf{Réf.} \\
    \midrule
    Laser ligne vert VLM-520-56 (520\,nm) & Commandé    & 77.08\,CHF & DigiKey \\
    2$\times$ caméra USB OV9726 (720p)\textsuperscript{*} & Commandé & 17.60\,CHF & Galaxus \\
    Adaptateur Gigabit Ethernet USB-C      & Commandé    & 8.90\,CHF  & Galaxus \\
    Adaptateur encastrable USB-C BU-BU     & Commandé    & 15.90\,CHF & Galaxus \\
    Adaptateur de panneau USB-A f-f        & Commandé    & 13.70\,CHF & Galaxus \\
    Contribution climatique                & Commandé    & 0.65\,CHF  & Galaxus \\
    \midrule
    Raspberry Pi 4 Model B                 & Stock       & $\approx$\,65\,CHF & --- \\
    Pi Camera Module 3                     & Stock       & $\approx$\,35\,CHF & --- \\
    Caméra USB Arducam IMX586              & Prêt        & $\approx$\,50\,CHF & --- \\
    Raspberry Pi Touch Display 2           & Stock       & $\approx$\,70\,CHF & --- \\
    Moteur pas à pas NEMA 17 (17HS3401)    & Stock       & $\approx$\,12\,CHF & --- \\
    Driver moteur DM320T                   & Stock       & $\approx$\,45\,CHF & --- \\
    Alimentation 24\,V 24\,W (Hama)        & Stock       & $\approx$\,25\,CHF & --- \\
    Convertisseur 24$\to$5\,V LM2596       & Stock       & $\approx$\,4\,CHF  & --- \\
    Interrupteur, connecteurs, capteur capot & Stock     & $\approx$\,15\,CHF & --- \\
    PCB dédié (fabrication)                & Stock       & $\approx$\,15\,CHF & --- \\
    Matière (PLA, panneau MDF, visserie)   & Stock       & $\approx$\,30\,CHF & --- \\
    \midrule
    \textbf{Total estimé (prix de marché)} & & \textbf{$\approx$\,500\,CHF} & \\
    \bottomrule
  \end{tabular}
  \caption{Nomenclature des pièces et estimation des prix. Les prix « Stock » sont estimés au prix de marché ; ils n'entrent pas dans le coût réel pour l'école.}
  \label{tab:budget-bom}
\end{table}

\textsuperscript{*}\,Les deux caméras USB OV9726 ont bien été commandées et
payées (elles figurent donc dans le coût réel), mais n'ont \emph{pas} pu être
utilisées : leur angle d'ouverture, non précisé à la commande, s'est avéré trop
étroit ($\approx 40^\circ$) pour couvrir la zone de scan. La seconde caméra du
montage final est une \textbf{Arducam IMX586} prêtée par l'école, qui n'a donc
engendré aucune dépense.

\subsection{Coût réel pour l'école}

Les heures de travail ne sont pas facturées et les moyens de fabrication du
MakeLab ne génèrent pas de dépense supplémentaire. Le coût réellement engagé
se limite aux deux commandes passées pour le projet, dont les quittances
figurent dans les annexes.

\begin{table}[H]
  \centering
  \begin{tabular}{@{}p{8cm}r@{}}
    \toprule
    \textbf{Commande} & \textbf{Montant (TTC)} \\
    \midrule
    DigiKey --- laser ligne VLM-520-56          & 77.08\,CHF \\
    Galaxus --- caméras USB et adaptateurs       & 56.75\,CHF \\
    \midrule
    \textbf{Coût réel total}                     & \textbf{133.83\,CHF} \\
    \bottomrule
  \end{tabular}
  \caption{Coût réel engagé pour le projet (quittances en annexe).}
  \label{tab:budget-reel}
\end{table}

Le coût réel de \textbf{133.83\,CHF} respecte largement l'objectif
budgétaire \textbf{O5} (\(\leq\) 300\,CHF). Même en valorisant l'ensemble
des composants au prix de marché ($\approx$\,500\,CHF,
\autoref{tab:budget-bom}), le coût matériel reste du même ordre de grandeur
que l'objectif, l'écart provenant essentiellement de l'écran tactile et du
Raspberry Pi.
